(1) 在TT规范的自由测试质量坐标近似不变，空间线元$dl^2=[1+h_+(t)]dx^2+[1-h_+(t)]dy^2$。两臂固有长度为$L_x=L\sqrt{1+h_+}\simeq L(1+h_+/2)$、$L_y=L\sqrt{1-h_+}\simeq L(1-h_+/2)$。所以差动长度$\Delta L=L_x-L_y\simeq Lh_+(t)=Lh_0\cos\Omega t$，公共尺度部分在一阶抵消。
(2) 在干涉仪中心的局部自由落体系，横向相对位移满足$\ddot\xi^i=\frac12\ddot h_{ij}^{TT}\xi_0^j$，其中$\xi_0$为未扰动分离。对沿$x$的基线有$\ddot{\delta\xi_x}=\frac12L\ddot h_+$，取无额外自由振荡的初始条件得$\delta\xi_x=Lh_+/2$；沿$y$得$\delta\xi_y=-Lh_+/2$，故差动长度同样为$Lh_+$。两个规范图像的可观测响应一致。
(3) 忽略有限光传播时间内$h_+$的变化时，长度变化$\delta L_i$使该臂往返光程改变$2\delta L_i$，往返相位$\delta\phi_i=(2\pi/\lambda)\,2\delta L_i$。因此干涉两臂差动相位$\Delta\phi=(4\pi/\lambda)(\delta L_x-\delta L_y)=(4\pi Lh_0/\lambda)\cos\Omega t$。该主阶相位随臂长增加，但实际信噪还受光学腔储能与散粒噪声影响。
(4) 若$\Omega L/c$不小，光在臂内往返时引力波已发生显著相位变化，需积分沿光线传播时的度规扰动，响应将包含频率相关的sinc型传递函数；不能直接以$Lh_+$代替整个相位过程。若探测器两臂相对偏振主轴旋转角$\alpha$，差动主阶应乘$\cos2\alpha$；加号偏振主轴偏离45度时灵敏度下降，交叉偏振还需另一几何投影。
